\documentclass[11pt]{article}
\usepackage[T1]{fontenc}
\usepackage{amsmath,amssymb,geometry}
\DeclareFontShape{OMX}{cmex}{m}{n}{<-> cmex10}{}
\usepackage[hidelinks]{hyperref}
\geometry{margin=1.05in}
\setlength{\emergencystretch}{3em}
\title{Eventual extension concavity of balanced arrays in fixed dimension}
\author{ProgramComputer}
\date{November 2026}
\begin{document}
\maketitle
\begin{abstract}
For each fixed number of columns, we prove strict concavity of the
effective extension counts at every genuine internal quota once the
number of rows is sufficiently large, uniformly over all proper caps.
The proof controls moving lower and upper quotas. Near the upper
boundary the exact falling factorial and the tilted entropy contribute
negative terms of different orders, which remain valid at their
crossover.
\end{abstract}
\section{The all-quota theorem}\label{sec:proof}

For nonnegative integer $n$, integers $d\ge2$ and $1\le r\le d$, let $a_k$ count
labeled binary $n$ by $d$ arrays with every column sum $k$ and every row
sum at most $r$. Write $D=\lfloor nr/d\rfloor$ and
$u_k=(k+1)(a_{k+1}/a_k)^{1/d}$ for $0\le k<D$.
The counts through $D$ are positive by distributing the successive columns' ones into consecutive blocks of a cyclic ordering of the rows.

\textbf{Theorem.} For each $d\ge2$ there exists $N(d)$ such that, simultaneously for every $1\le r<d$ and every $n\ge N(d)$,

\begin{equation*}
u_{k-1}-2u_k+u_{k+1}<0
\qquad (1\le k\le D-2).
\end{equation*}

Empty index ranges are vacuous. At full cap $r=d$ the same second differences are zero, since $a_k=\binom nk^d$ and $u_k=n-k$. A zero extension at quota $D$ is outside the statement. The theorem addresses an extension inequality for the original balanced-array coefficients. \href{https://www.symmetricfunctions.com/realRootedOpenProblems.htm#conjBalancedArrays}{Per Alexandersson's original real-rootedness question}

\subsection{Rank one and compact interior}

If $r=1$, the extension count raised to the $d$th power is
$\prod_{j=0}^{d-1}(n-dk-j)$. On any positive interval containing three genuine extensions, the geometric mean of those affine factors is strictly concave. In fact its second derivative divided by the mean is $d^2$ times the negative variance of the $d$ reciprocals of the factors under their uniform distribution. The reciprocals are distinct since $d\ge2$.

Now fix $2\le r<d$ and put $\rho=r/d$. The \href{https://github.com/ProgramComputer/open-math-research/blob/27e3e70d6603e8a395c5fdb7ff944922bce83380/topics/balanced-arrays/engset-curvature/finite-extension.md#fixed-dimensions-compact-interior}{compact-interior extension argument} proves strictness uniformly on every compact subinterval of $0<k/n<\rho$, for all sufficiently large $n$. The tilted entropy there gives $u_k=n f(k/n)+f_1(k/n)+n^{-1} f_2(k/n)+O(n^{-2})$, with smooth retained terms and $f''<0$. The two neighborhoods with distance to the nearest quota boundary tending to infinity are handled by the \hyperref[sec:lower]{lower mesoscopic proof} and the \hyperref[sec:upper]{upper mesoscopic proof}.

\subsection{Fixed lower quota}

Put $p_k=a_k/\binom nk^d$. Independently select a uniform $k$-subset of rows in each column. At fixed $k$, a specified row violates the cap with probability $\binom d{r+1}(k/n)^{r+1}+O(n^{-r-2})$. For two specified rows the joint probability is $O(n^{-2(r+1)})$. To see this, prescribe $r+1$ occupied columns for each row in finitely many ways. In each column, requiring one or both specified labels has probability $O(n^{-1})$ or $O(n^{-2})$ respectively. First and second Bonferroni bounds and the logarithm give

\begin{equation*}
\log p_k=-\binom d{r+1} k^{r+1}n^{-r}+O(n^{-r-1}).
\end{equation*}

Since $u_k=(n-k)(p_{k+1}/p_k)^{1/d}$, the second difference at fixed $k\ge1$ equals

\begin{equation*}
-\frac1d\binom d{r+1} n^{1-r}
\left.\Delta^3 t^{r+1}\right|_{t=k-1}+O(n^{-r}),
\end{equation*}

where $\Delta f(t)=f(t+1)-f(t)$. The third difference is positive, since it is the integral of the strictly positive third derivative at positive arguments over a unit cube.

\subsection{Fixed upper deficit}

Put $e=nr-dk$, $b=\binom d{r-1}/\binom dr$, and
$\lambda=r(d-r)/(d(d-1))$. At any fixed feasible $e\ge0$ the main arrays before balancing use $e$ rows of size $r-1$ and all other rows of size $r$. The weights of all profiles with fewer deficient rows contribute an $O(1/n)$ relative factor after the balancing estimates below. Condition on the occupied sets of the deficient rows. For the remaining rows choose independent uniform size-$r$ subsets. Their projection to the first $d-1$ coordinate counts has covariance $\lambda(I-\mathbf1\mathbf1^{\mathrm T}/d)$, of determinant $\lambda^{d-1}/d$. Its support differences include every coordinate unit vector by exchanging a label with label $d$. Thus its characteristic function has a single modulus-one point on the torus. Near that point the covariance is nonsingular. Fourier inversion and Gaussian scaling give, for every target at bounded displacement from its expectation, probability asymptotic to $C_{d,r}n^{-(d-1)/2}$, with $C_{d,r}=\sqrt d(2\pi\lambda)^{-(d-1)/2}$. The same argument gives a uniform upper bound of this order, even for the other finitely many deficit profiles. The last column equation follows from the total sum. Therefore

\begin{equation*}
a_k=C_{d,r}\binom{d}{r}^{n} n^{e-(d-1)/2}\frac{b^e}{e!}(1+o(1)).
\end{equation*}

At each fixed $D-k\ge2$, pass if necessary to a fixed residue of $nr\bmod d$. Then $u_k$ tends to
$(\rho/b)\bigl(e(e-1)\cdots(e-d+1)\bigr)^{1/d}$, and the neighboring extensions tend to the same expression at $e\pm d$. Those three deficits are at least $d$. The limiting function is strictly concave on this interval by the same geometric-mean derivative calculation. There are only finitely many possible residues.

\subsection{Moving upper deficit and the crossover}

For $e\to\infty$ with $e/n\to0$, the \hyperref[sec:upper]{upper-edge expansion} gives

\begin{equation*}
u_{k-1}-2u_k+u_{k+1}=
-\frac{2d(r-1)(d-r)(d-r+1)}{r(d-r+2)n}
-\frac{d^2(d^2-1)\rho}{12b e^3}
+o(n^{-1}+e^{-3}).
\end{equation*}

The first term comes from the curvature of the tilted entropy. The second comes from preserving the exact falling factorial in the adjacent ratio. If $e/n^{1/3}\to t\in(0,\infty)$, the coefficient of $n^{-1}$ is the sum of the first negative constant and the second negative constant multiplied by $t^{-3}$. Below the crossover the factorial term dominates. Above it the entropy term dominates.

\subsection{Combining the regimes}

If the claimed row threshold failed for a fixed $d$, there would be a sequence of failures with $n\to\infty$. There are finitely many proper caps, so restrict to a subsequence with fixed $r$. Rank one is already handled. The remaining densities $k/n$ lie in $[0,\rho]$, so along a subsequence they converge. An interior limit lies in one compact interval, contradicting compact-interior strictness. A limit of zero has a further subsequence with constant $k$ or with $k\to\infty$. These are the fixed lower quota and \hyperref[sec:lower]{lower mesoscopic} cases. For a limit of $\rho$, the integer deficit $nr-dk$ is constant along a further subsequence or diverges to infinity. The finite residue argument and the moving upper expansion cover those cases. Each possible failure leads to a contradiction.

This proof gives no effective $N(d)$ and no uniform cutoff over dimensions. The all-parameter inequality $2u_k\ge u_{k-1}+u_{k+1}$ for every finite admissible $(n,r,d,k)$ remains open. These arguments also leave the all-cap real-rootedness conjecture open.

\section{Finite extension concavity at the lower mesoscopic edge}\label{sec:lower}

\textbf{Claim.} Fix integers $2\le r<d$. For every sequence $k\to\infty$ with $k/n\to0$, the effective extension counts $u_k=(k+1)(a_{k+1}/a_k)^{1/d}$ are eventually strictly concave at $k$. With $x=k/n$, more precisely,

\begin{equation*}
u_{k+1}-2u_k+u_{k-1}
=-\binom{d-1}{r}r(r-1)\frac{x^{r-2}}{n}
\left(1+O_{d,r}\left(x+\frac1k\right)\right).
\end{equation*}

For $r=1$, the separate geometric-mean argument applies. This note handles simultaneous growth of $k$ and $n/k$ at fixed dimension and cap.

\subsection{A parameter that isolates the conditioning correction}

Give every subset $A\subseteq[d]$ row weight $1$ if $|A|\le r$ and $s$ otherwise, where $0\le s\le1$. Define

\begin{equation*}
F_s(z_1,\ldots,z_d)=\sum_A \omega_s(A)\prod_{i\in A}z_i, \qquad
G_s(z)=F_s(z,\ldots,z).
\end{equation*}

With all variables tilted by $z$, the resulting row vector $X$ has common coordinate mean

\begin{equation*}
x=\frac{zG_s'(z)}{dG_s(z)}.
\end{equation*}

The inverse tilt $z=z_s(x)$ is analytic in $x$ on a fixed small disk, uniformly over $s\in[0,1]$, with $z_s(x)=x+O(x^2)$. This follows from the analytic implicit function theorem since the map $z\mapsto zG_s'(z)/(dG_s(z))$ has derivative one at zero for each $s$, and the parameter interval is compact. Differentiation in $s$ shows $\partial_s z_s(x)=O(x^{r+1})$ because the weights first change at degree $r+1$.

For real $x>0$, set

\begin{equation*}
K_s(x,\theta)=
\log F_s(z_s(x)e^{i\theta_1},\ldots,z_s(x)e^{i\theta_d})
-\log G_s(z_s(x))-ix\sum_i\theta_i .
\end{equation*}

All logarithms can be defined through the convergent local series at one, uniformly for $\theta\in[-\pi,\pi]^d$ and $x$ in a fixed small complex disk. The constant and linear $\theta$-terms vanish exactly, the latter because the tilt is retuned to mean $x$. The same is true after differentiating in $s$ and $x$. The analytic expressions $\partial_\theta^2 K_s$ vanish to order at least one in $x$, and $\partial_s\partial_\theta^2K_s$ vanish to order at least $r+1$. For the latter, differentiate $F_s$ directly and through the tilt, using that the first changed weights are at degree $r+1$ and $\partial_s z_s=O(x^{r+1})$. Taylor in $\theta$ along the segment from zero, and analytic bounds for the $x$-derivatives, give for $0\le j\le3$

\begin{equation*}
|\partial_x^jK_s|\le C_j x^{1-j}\lVert \theta\rVert ^2, \qquad
|\partial_s\partial_x^j K_s|\le C_j x^{r+1-j}\lVert \theta\rVert ^2 .
\end{equation*}

These bounds hold all the way across the torus cube, with a constant depending on fixed $d,r$.

The row vector has probability bounded below of being zero and probability at least $c x$ of taking each unit vector. Comparing an independent pair of such vectors gives

\begin{equation*}
1-|\mathbb E e^{i\theta\cdot X}|^2
\ \ge\ c x\sum_i(1-\cos\theta_i)\ \ge\ c' x\lVert \theta\rVert ^2 .
\end{equation*}

Therefore $\Re K_s\le-c''x\lVert \theta\rVert ^2$ throughout the cube. Set

\begin{equation*}
I_s(x)=(2\pi)^{-d}\int_{[-\pi,\pi]^d}e^{nK_s(x,\theta)}\thinspace d\theta
\end{equation*}

for real $x$ without restricting it to lattice points. This integral is real by conjugation. Uniformly for $nx$ sufficiently large, $I_s(x)\ge C(nx)^{-d/2}$. To see this, change variables to $v=\sqrt{nx}\thinspace \theta$. Near zero the centered exponent equals $-v^\mathrm{T}(\mathrm{Cov}(X)/x)v/2$ plus $O(\lVert v\rVert ^3/\sqrt{nx})$. Analytic boundedness of third $\theta$-derivatives divided by $x$ gives the error, and $\mathrm{Cov}(X)/x=I_d+O(x)$ from the singleton masses. Dominated Gaussian integration uniformly over the interval of $s$ and sufficiently small $x$ shows that the scaled integral differs by $o(1)$ from the positive Gaussian integral with that covariance. This also supplies positivity for off-grid $x$.

\subsection{Bounds on derivatives of the log-conditioning ratio}

Differentiate the integral before restricting to the lattice. For $0\le j\le3$, each product generated by differentiating $e^{nK_s}$ has polynomial factors in $nx\lVert \theta\rVert ^2$ times the Gaussian bound and an outside factor $x^{-j}$. If one $s$-derivative is taken, exactly one factor carries an extra $x^r$. Integrating the bounds gives

\begin{equation*}
|\partial_x^jI_s(x)|\le C_j x^{-j}(nx)^{-d/2}, \qquad
|\partial_s\partial_x^j I_s(x)|\le C_j x^{r-j}(nx)^{-d/2}.
\end{equation*}

Together with the positive lower bound, quotient differentiation of $\partial_s\log I_s$ and integration over $s$ yield

\begin{equation*}
E_n(x):=\log I_0(x)-\log I_1(x), \qquad
|E_n^{(j)}(x)|\le C_j x^{r-j}\quad (0\le j\le3).
\end{equation*}

In particular the correction retains its small factor $x^r$. Comparing two independently approximated local limit errors would not give that factor.

At $x=k/n$, Fourier inversion gives the exact probability of every column sum being $k$, so with $\Phi_s(x)=\log G_s(z_s(x))-dx\log z_s(x)$,

\begin{equation*}
\log p_k=n(\Phi_0-\Phi_1)(x)+E_n(x), \qquad p_k=\frac{a_k}{\binom nk^d}.
\end{equation*}

The entropy derivative is $(\Phi_0-\Phi_1)'=-dg$, where $g(x)=\log(z_0(x)(1-x)/x)$ is the correction in \href{https://github.com/ProgramComputer/open-math-research/blob/27e3e70d6603e8a395c5fdb7ff944922bce83380/topics/balanced-arrays/engset-curvature/carried-load.md}{carried-load.md}. The Engset formula there gives

\begin{equation*}
g(x)=\binom{d-1}{r}x^r+O(x^{r+1})
\end{equation*}

as an analytic expansion at zero.

\subsection{Extracting the finite difference}

Let $h=1/n$. Extend $\log[p_{k+1}/p_k]/d$ as a smooth function of $x$ using the displayed exact representation at grid points. It equals $-g(x)+R_n(x)$, where

\begin{equation*}
R_n(x)=-n\int_0^h\bigl(g(x+t)-g(x)\bigr)\thinspace dt
+\frac1d\int_0^h E_n'(x+t)\thinspace dt.
\end{equation*}

For $0\le j\le2$ the bounds above give $R_n^{(j)}=O(hx^{r-1-j})$, uniformly when $nx$ is sufficiently large and $x$ sufficiently small. Thus the smooth extension of $u_k=(n-k)(p_{k+1}/p_k)^{1/d}$ is $n f(x)+Q_n(x)$, with

\begin{equation*}
f(x)=(1-x)e^{-g(x)}, \qquad
Q_n(x)=n f(x)(e^{R_n(x)}-1), \qquad Q_n''(x)=O(x^{r-3}).
\end{equation*}

Analyticity gives $f''(x)=-\binom{d-1}{r}r(r-1)x^{r-2}(1+O(x))$. Its symmetric finite difference is the integral of $f''$ against the triangular kernel of mass $h^2$ supported on $[-h,h]$. Over that window the leading power changes by a relative $O(h/x)$. Likewise the error from $Q_n$ is $O(h^2x^{r-3})$. Combining them gives the stated expansion and strict concavity along the entire lower mesoscopic edge.

The same derivative estimates also give the third forward difference of $\log p_k$ at this scale. On three consecutive steps, the leading contribution is $n h^3(\Phi_0-\Phi_1)'''(x)=-d\binom{d-1}{r}r(r-1)x^{r-2}/n^2$, with relative error $O(x+1/k)$, since the remainder has third derivative $O(x^{r-3})$. The extension inequality above does not rely on interpreting a discrete derivative at a noninteger quota.

\section{Strict extension concavity at the mesoscopic upper edge}\label{sec:upper}

Fix integers $2\le r<d$. Let $a_k$ count the labeled capped arrays, and let
$u_k=(k+1)(a_{k+1}/a_k)^{1/d}$. Put $e=nr-dk$, $h=e/n$, and suppose
$e\to\infty$ and $h\to0$. Then

\begin{equation*}
u_{k-1}-2u_k+u_{k+1}
=\frac{d^2H''(0)}n-\frac{d^2(d^2-1)\rho}{12b e^3}
+o(n^{-1}+e^{-3}), \qquad
\rho=\frac rd,\quad b=\frac{\binom d{r-1}}{\binom dr},
\end{equation*}

where $H''(0)<0$. Thus the proposed strict concavity holds along every such sequence, including the crossover $e\asymp n^{1/3}$. The current argument concerns fixed $d,r$.

\subsection{The moving tilted row law}

Let $P(y)=\sum_{j=0}^r\binom d{r-j}y^j/\binom dr$. A row with occupied set $S$ has tilted probability $y^{r-|S|}/(\binom dr P(y))$. Write $T=r-|S|$ for its deficit. Determine $y=y(h)$ by $yP'(y)/P(y)=h$. The implicit function theorem gives a real analytic $y$ near zero with $y(0)=0$, $y'(0)=1/b$. For positive small $h$, the size-$r$ subsets have individual probabilities bounded below, and each size-$(r-1)$ subset has probability bounded below by a positive constant times $h$. Also $\mathbb P(T=1)=h+O(h^2)$ and $\mathbb P(T\ge2)=O(h^2)$, with analytic probabilities.

Let $X$ be the row indicator vector, $Z$ its orthogonal projection to the sum-zero subspace in fixed orthonormal coordinates, and $x=(r-h)/d$. Then $\mathbb E X=x\mathbf1$, $\mathbb EZ=0$, and $\mathrm{Cov}(T,Z)=0$ by permutation invariance. The covariance of $Z$ is analytic in $h$ and positive definite at zero: uniform size-$r$ subsets have all the exchange differences $e_i-e_j$. Moreover $\mathrm{Var}(T)/h$ is analytic at zero with value one. For every $m\ge2$, the order-$m$ cumulant of $T$ equals $h+O(h^2)$, since $\log\mathbb E e^{vT}=h(e^v-1)+O(h^2)$ as a convergent analytic function near $v=0$. Every joint cumulant involving at least one copy of $T$ vanishes at $h=0$. Constants here and below may depend on the fixed $d,r$.

\subsection{A two-scale expansion with analytic coefficients}

Write $K_h(\theta)=\mathbb E e^{i\theta\cdot X}$ for the uncentered characteristic function. Its squared modulus satisfies

\begin{equation*}
1-|K_h(\theta)|^2
=\sum_{S,S'}\mathbb P(S)\mathbb P(S')
\bigl(1-\cos(\theta\cdot(\mathbf1_S-\mathbf1_{S'}))\bigr).
\end{equation*}

For each $j<d$ a pair of size-$r$ subsets differing by exchange of $j$ and $d$ gives a fixed positive multiple of $1-\cos(\theta_j-\theta_d)$. A size-$r$ set containing $d$ and the set obtained by deleting $d$ give a positive multiple of $h(1-\cos\theta_d)$. All terms are nonnegative, so, taking the relevant torus representatives in $[-\pi,\pi]$, these inequalities imply

\begin{equation*}
|K_h(\theta)|^n
\le \exp\left(-c n\sum_{j<d}|\theta_j-\theta_d|_{2\pi}^2
-c e|\theta_d|_{2\pi}^2\right)
\end{equation*}

for some $c>0$. The coordinates $\theta_j-\theta_d$ and $\theta_d$ preserve torus measure. In particular the bound retains the transverse volume factor $n^{-(d-1)/2}$ even when $e$ grows arbitrarily slowly compared with $n$. Near the origin a constant invertible linear change gives the argument $(v,w)$ for the random vector $(T-h,Z)$, where $v$ is minus the arithmetic mean of the angles. The quadratic bound in torus coordinates dominates a positive constant times $e v^2+n\lVert w\rVert ^2$ uniformly for $0<h\le h_0$. This follows since each transverse angle difference controls the entire sum-zero projection, and the difference of $v$ from $-\theta_d$ is a fixed linear function of that projection with $e\le n$.

Let $I(n,e)$ be the tilted probability of column sum $(k,\ldots,k)$. Fourier inversion in $d$ dimensions is centered exactly at the mean. On the central box substitute $v=t/\sqrt e$, $w=s/\sqrt n$. The Jacobian is a positive constant times $e^{-1/2}n^{-(d-1)/2}$. Outside the box $|t|+\lVert s\rVert \le e^{1/100}$, the torus bound above, integrated using the two Gaussian scales, contributes at most that same volume factor times $O(\exp(-c'e^{1/50}))$. To justify use of this box globally, first restrict to a fixed neighborhood of the torus origin where the representatives do not wrap. On its complement one angular difference or the common angle is bounded away from zero, so the same product estimate gives the two-scale volume times $O(e^{-c'' e})$.

The centered row log-characteristic function is jointly analytic in a fixed complex neighborhood of $(h,v,w)=0$, using its branch at zero. A term of degree $a$ in $v$ and $b'$ in $w$, of total degree $\ell=a+b'\ge2$, has after multiplication by $n$ and rescaling the form

\begin{equation*}
e^{1-\ell/2}\begin{cases}
h^{b'/2}\times\text{an analytic coefficient in }h,& a\ge1,\\
h^{b'/2-1}\times\text{an analytic coefficient in }h,&a=0.
\end{cases}
\end{equation*}

The first case uses the vanishing of each radial cumulant at $h=0$. In the second case $b'\ge2$. At degree two the scaled Gaussian is block diagonal, with blocks $\mathrm{Var}(T)/h$ and $\mathrm{Cov}(Z)$, both bounded and positive definite after restricting $h_0$. Expand through degree eleven. The Taylor remainder, after multiplication by $n$, is $O(e^{-5}(|t|+\lVert s\rVert )^{12})$ on the central box. For this bound, separate the log-characteristic function at $v=0$ and its difference from that restriction. The former remainder has degree at least twelve in $w$, while the latter is divisible by $h$ and its remainder is bounded by a constant times $h(|v|+\lVert w\rVert )^{12}$. Uniform convergence of the analytic Taylor series gives these bounds.

Now expand the exponential of the nonquadratic terms through order nine in $e^{-1/2}$. The omitted terms, after multiplication by the Gaussian, have integrated size $O(e^{-5})$: on the box the nonquadratic terms are uniformly $o(1)$, and the remainder is bounded by $e^{-5}$ times a fixed polynomial in $|t|+\lVert s\rVert $ times the Gaussian. Extend each polynomial Gaussian integral to all of $\mathbb R^d$ with exponentially small error. Odd powers of $e^{-1/2}$ vanish by simultaneous reflection, since their total monomial degree is odd. Terms of odd total transverse degree vanish by transverse reflection alone because the Gaussian is block diagonal. The surviving powers of $h$ displayed above are therefore nonnegative integers. The Gaussian integrals, including the leading determinant factor, are analytic in $h$, since their covariance matrices are analytic and nonsingular at zero. The leading coefficient is positive on a neighborhood of zero.

Taking the logarithm gives analytic functions $f_j$ such that

\begin{equation*}
\log I(n,e)=-\tfrac12\log e-\tfrac{d-1}2\log n
+\sum_{j=0}^4 f_j(h)e^{-j}+O(e^{-5}).
\end{equation*}

The error is uniform for $0<h\le h_0$ and $e\to\infty$. At $h=0$ all nonquadratic terms with transverse variables vanish. The pure radial normalized cumulants tend to one. In the leading Gaussian law $t$ is standard normal, so the relative coefficient of $e^{-1}$ is $\mathbb E[t^4]/24-\mathbb E[t^6]/72=3/24-15/72=-1/12$. Taking the logarithm preserves that first relative coefficient. Thus $f_1(0)=-1/12$.

\subsection{Isolating the factorial and taking second differences}

Write $S(h)=\log P(y)-h\log y$. The exact tilted counting identity is
$a_k=\binom{d}{r}^{n}\exp(n S(h))I(n,e)$. Since $S'(h)=-\log y$, one has
$S(h)=h(1+\log(b/h))+\Psi(h)$ with $\Psi$ analytic, $\Psi(0)=0$, and
$\Psi'(h)=-\log(by(h)/h)$, so $\Psi'(0)=0$. Stirling's formula to order four, with error $O(e^{-5})$, now gives

\begin{equation*}
\log a_k=n\log\binom dr+e\log(nb)-\log(e!)
-\tfrac{d-1}2\log n+n\Psi(h)
+\sum_{j=0}^4 F_j(h)e^{-j}+O(e^{-5}),
\end{equation*}

where the $F_j$ are analytic and $F_1(0)=0$ by cancellation of the $-1/(12e)$ term.

Put $\epsilon=1/e$. Replacing $k$ by $k+1$ replaces $e$ by $e-d$ and $h$ by $h'=h(1-d\epsilon)$. After taking the adjacent ratio, the factorial contribution to $u_k/e$ is $\prod_{j=0}^{d-1}(1-j\epsilon)^{1/d}$. The remaining prefactor is $(x+h\epsilon)/b$. The contribution in the logarithm from $n\Psi$ after taking the $d$th root is

\begin{equation*}
\frac{\Psi(h-dh\epsilon)-\Psi(h)}{d h\epsilon}
=-\int_0^1 \Psi'(h-dh\epsilon t)\thinspace dt,
\end{equation*}

which extends analytically at $h\epsilon=0$. The other displayed log terms give the analytic contributions $\epsilon^j[(1-d\epsilon)^{-j}F_j(h')-F_j(h)]/d$. Thus $u_k/e$ is an analytic function of $(h,\epsilon)$ up to $O(e^{-5})$. Its expansion in $\epsilon$ yields analytic $B,C,E_2,E_3$ with

\begin{equation*}
u_k=nH(h)+B(h)+C(h)e^{-1}+E_2(h)e^{-2}+E_3(h)e^{-3}+O(e^{-4}),
\qquad H(h)=xy(h).
\end{equation*}

Indeed the constant coefficient of $u_k/e$ is $(x/b)e^{-\Psi'(h)}=xy(h)/h$. At $h=0$ the $\Psi$ contribution vanishes, and so do the differences for $j=0,1$, by $h'=h=0$ and $F_1(0)=0$. For $j\ge2$ the difference starts at order $\epsilon^{j+1}$. Therefore $C(0)$ is $\rho/b$ times the second coefficient in $\prod_{j=0}^{d-1}(1-j\epsilon)^{1/d}$. Expanding its logarithm gives

\begin{equation*}
C(0)=-\frac{(d^2-1)\rho}{24b}.
\end{equation*}

For the other constant, write $s_{\ast}=d-r$ and $b_2=\binom d{r-2}/\binom dr$. The deficit equation has expansion $h=b y+(2b_2-b^2)y^2+O(y^3)$. Inverting gives $y=h/b-(2b_2-b^2)h^2/b^3+O(h^3)$. Multiplying by $x=\rho-h/d$ gives

\begin{equation*}
H''(0)
=\frac{2((r-1)b^2-2r b_2)}{d b^3}
=-\frac{2(r-1)s_{\ast}(s_{\ast}+1)}{d r(s_{\ast}+2)}<0.
\end{equation*}

Finally the centered second difference of step $d$ in $e$ of $nH(e/n)$ is $d^2H''(0)/n+o(n^{-1})$. For $B(e/n)$ it is $O(n^{-2})$. For $C(e/n)/e$ its second derivative is $2C(0)/e^3+o(e^{-3})$ uniformly on $[e-d,e+d]$, so its second difference contributes $2d^2 C(0)/e^3+o(e^{-3})$. The second derivatives of the next two smooth terms are $O(e^{-4})$, using $h\to0$. The pointwise uniform remainder also has second difference $O(e^{-4})$. This proves the asserted two-scale asymptotic.

\paragraph{Sources.}
The original counting polynomial is due to Per Alexandersson in his
\href{https://www.symmetricfunctions.com/realRootedOpenProblems.htm\#conjBalancedArrays}{uniform-matroid conjecture}.
The interior argument uses the carried-load correction discussed in the
linked analysis, with Engset background in P. Azimzadeh and T. Carpenter,
\emph{Fast Engset computation},
\href{https://doi.org/10.1016/j.orl.2016.02.011}{doi:10.1016/j.orl.2016.02.011}.
\end{document}
